\documentclass[a4paper,11pt]{article}

\usepackage[a4paper,margin=2cm]{geometry}
\usepackage[T1]{fontenc}
\usepackage[utf8]{inputenc}
\usepackage[ngerman,english]{babel}
\usepackage{lmodern}
\usepackage{microtype}
\usepackage{xcolor}
\usepackage{parskip}
\usepackage{enumitem}
\usepackage[most]{tcolorbox}
\usepackage{titlesec}
\usepackage{array}
\usepackage{longtable}
\usepackage[table]{xcolor}
\usepackage[hidelinks]{hyperref}

\definecolor{HeaderBlueDark}{RGB}{70,110,170}
\definecolor{RowBlueLight}{RGB}{235,243,255}
\definecolor{RowBlueDark}{RGB}{220,233,250}
\definecolor{SoftGray}{RGB}{90,90,90}

\setlength{\parindent}{0pt}
\setlist[itemize]{leftmargin=1.4em,itemsep=0.35em,topsep=0.35em}
\linespread{1.06}
\renewcommand{\arraystretch}{1.35}
\setlength{\tabcolsep}{5pt}

\titleformat{\section}[block]
{\Large\bfseries\color{HeaderBlueDark}}
{}{0pt}{}
\titlespacing{\section}{0pt}{1.3em}{0.8em}

\titleformat{\subsection}[block]
{\large\bfseries\color{HeaderBlueDark}}
{}{0pt}{}
\titlespacing{\subsection}{0pt}{1.0em}{0.6em}

\newtcolorbox{exbox}{enhanced,breakable,colback=RowBlueLight,colframe=RowBlueDark,boxrule=0.4pt,arc=1.5mm,left=1.2mm,right=1.2mm,top=1mm,bottom=1mm,title={Beispiele \textnormal{(Examples)}},fonttitle=\bfseries,colbacktitle=RowBlueDark,coltitle=black}

\NewDocumentEnvironment{grammarentry}{mm+b}{%
    \begin{tcolorbox}[enhanced,breakable,colback=white,colframe=HeaderBlueDark,boxrule=0.6pt,arc=1.5mm,left=1.5mm,right=1.5mm,top=1mm,bottom=1mm,colbacktitle=HeaderBlueDark,coltitle=white,fonttitle=\bfseries,title={#1\hfill\normalfont\footnotesize #2}]
    #3
    \end{tcolorbox}
}{}

\newcommand{\GermanText}[1]{\textbf{Deutsch: }#1\par}
\newcommand{\EnglishText}[1]{{\small\color{SoftGray}\textbf{English: }#1\par}}
\newcommand{\ExampleItem}[2]{\item \textbf{#1}\\{\small\color{SoftGray}#2}}

\title{Nomenendungen im Deutschen\\\large \textnormal{Noun Endings in German}}
\date{}

\begin{document}
\maketitle
\tableofcontents
\newpage

\section{Grundidee \textnormal{(Basic idea)}}
\begin{grammarentry}{Was verändert sich am Nomen?}{Kasusbedingte Nomenendungen}
\GermanText{Im Deutschen verändert sich das Nomen selbst wegen des Kasus nur in bestimmten Gruppen. Die wichtigsten Fälle sind: \textbf{Genitiv Singular mit -s/-es}, \textbf{Dativ Plural mit -n}, \textbf{gemischte Deklination} und einige \textbf{Sonderfälle}.}
\EnglishText{In German, the noun itself changes because of case only in certain groups. The most important cases are: \textbf{genitive singular with -s/-es}, \textbf{dative plural with -n}, \textbf{mixed declension}, and a few \textbf{special cases}.}

\GermanText{Die Endungen von Artikeln und Adjektiven sind ein anderes Thema. In \glqq mit den kleinen Kindern\grqq sind \textbf{-en} an \glqq kleinen\grqq und \textbf{-n} an \glqq Kindern\grqq nicht dieselbe grammatische Erscheinung.}
\EnglishText{Article and adjective endings are a different topic. In \emph{mit den kleinen Kindern}, the \textbf{-en} on \emph{kleinen} and the \textbf{-n} on \emph{Kindern} are not the same grammatical phenomenon.}
\end{grammarentry}

\begin{grammarentry}{Wichtig}{Pluralbildung ist nicht dasselbe}
\GermanText{Die normalen Pluralformen wie \textbf{Tag -- Tage}, \textbf{Kind -- Kinder}, \textbf{Frau -- Frauen} oder \textbf{Auto -- Autos} muss man im Allgemeinen mit dem Nomen lernen. Sie sind keine normalen Kasusendungen. Eine wichtige Ausnahme ist das zusätzliche \textbf{-n im Dativ Plural}.}
\EnglishText{Normal plural forms such as \textbf{Tag -- Tage}, \textbf{Kind -- Kinder}, \textbf{Frau -- Frauen}, or \textbf{Auto -- Autos} generally have to be learned with the noun. They are not ordinary case endings. One important exception is the additional \textbf{-n in the dative plural}.}
\end{grammarentry}

\section{Gesamtüberblick \textnormal{(Overall overview)}}
\rowcolors{2}{RowBlueLight}{RowBlueDark}
\begin{longtable}{>{\bfseries}p{3.1cm}|p{3.1cm}|p{3.2cm}|p{5.5cm}}
\rowcolor{HeaderBlueDark}
\color{white}\textbf{Typ (Type)} &
\color{white}\textbf{Kasus (Case)} &
\color{white}\textbf{Endung (Ending)} &
\color{white}\textbf{Beispiel (Example)} \\
Starke mask./neutr. Nomen & Genitiv Singular & -s oder -es & der Mann $\rightarrow$ des Mannes; das Auto $\rightarrow$ des Autos \\
Dativ Plural & Dativ Plural & meist zusätzlich -n & die Kinder $\rightarrow$ mit den Kindern \\
Gemischte Deklination & Genitiv zusätzlich -s & -ns / -ens & der Name $\rightarrow$ des Namens \\
Sonderfall \glqq Herz\grqq & Dat./Gen. Singular & -en / -ens & das Herz $\rightarrow$ dem Herzen, des Herzens \\
Eigennamen & Genitiv & meist -s & Annas Buch; Peters Auto \\
\end{longtable}

\section{Genitiv Singular: -s und -es \textnormal{(Genitive singular: -s and -es)}}
\begin{grammarentry}{Grundregel}{Maskulin und Neutrum}
\GermanText{Viele maskuline und neutrale Nomen bekommen im Genitiv Singular \textbf{-s} oder \textbf{-es}. Feminine Nomen bekommen normalerweise keine Genitivendung am Nomen.}
\EnglishText{Many masculine and neuter nouns take \textbf{-s} or \textbf{-es} in the genitive singular. Feminine nouns normally do not take a genitive ending on the noun itself.}
\end{grammarentry}

\rowcolors{2}{RowBlueLight}{RowBlueDark}
\begin{longtable}{>{\bfseries}p{3cm}|p{3.3cm}|p{3.3cm}|p{5.3cm}}
\rowcolor{HeaderBlueDark}
\color{white}\textbf{Genus (Gender)} &
\color{white}\textbf{Nominativ (Nominative)} &
\color{white}\textbf{Genitiv (Genitive)} &
\color{white}\textbf{Bemerkung (Remark)} \\
Maskulin & der Mann & des Mannes & -es \\
Maskulin & der Lehrer & des Lehrers & -s \\
Neutrum & das Kind & des Kindes & -es \\
Neutrum & das Auto & des Autos & -s \\
Feminin & die Frau & der Frau & keine Endung am Nomen \\
\end{longtable}

\begin{grammarentry}{Wann -es?}{Praktische Regel}
\GermanText{Bei vielen \textbf{einsilbigen} Nomen ist \textbf{-es} sehr häufig: \textbf{des Mannes, des Kindes, des Tages, des Buches}. Bei längeren Nomen ist oft \textbf{-s} üblich: \textbf{des Lehrers, des Computers, des Hotels}.}
\EnglishText{With many \textbf{one-syllable} nouns, \textbf{-es} is very common: \emph{des Mannes, des Kindes, des Tages, des Buches}. With longer nouns, \textbf{-s} is often usual: \emph{des Lehrers, des Computers, des Hotels}.}

\GermanText{Nach bestimmten Lauten und Schreibungen ist \textbf{-es} besonders wichtig oder deutlich natürlicher, zum Beispiel bei vielen Wörtern auf \textbf{-s, -ß, -x, -z, -tz, -sch}: \textbf{des Hauses, des Fußes, des Reflexes, des Satzes, des Tisches}.}
\EnglishText{After certain sounds and spellings, \textbf{-es} is especially important or clearly more natural, for example with many words ending in \textbf{-s, -ß, -x, -z, -tz, -sch}: \emph{des Hauses, des Fußes, des Reflexes, des Satzes, des Tisches}.}
\end{grammarentry}

\begin{exbox}
\begin{itemize}
\ExampleItem{Die Farbe des Autos ist schön.}{The color of the car is beautiful.}
\ExampleItem{Das Ende des Tages war ruhig.}{The end of the day was quiet.}
\ExampleItem{Die Meinung der Frau ist wichtig.}{The woman's opinion is important.}
\end{itemize}
\end{exbox}

\section{Gemischte Deklination: Name, Gedanke usw. \textnormal{(Mixed declension: Name, Gedanke, etc.)}}
\begin{grammarentry}{Besondere maskuline Nomen}{-n/-en plus Genitiv -s}
\GermanText{Einige maskuline Nomen bekommen im \textbf{Akkusativ und Dativ Singular -n/-en} und im \textbf{Genitiv Singular zusätzlich -s}. Typische Wörter sind \textbf{der Name, der Gedanke, der Glaube, der Wille, der Buchstabe, der Same}.}
\EnglishText{Some masculine nouns take \textbf{-n/-en in the accusative and dative singular} and an \textbf{additional -s in the genitive singular}. Typical words are \emph{der Name, der Gedanke, der Glaube, der Wille, der Buchstabe, der Same}.}
\end{grammarentry}

\rowcolors{2}{RowBlueLight}{RowBlueDark}
\begin{longtable}{>{\bfseries}p{3.2cm}|p{5cm}|p{5cm}}
\rowcolor{HeaderBlueDark}
\color{white}\textbf{Kasus (Case)} &
\color{white}\textbf{der Name} &
\color{white}\textbf{der Gedanke} \\
Nominativ & der Name & der Gedanke \\
Akkusativ & den Namen & den Gedanken \\
Dativ & dem Namen & dem Gedanken \\
Genitiv & des Namens & des Gedankens \\
\end{longtable}

\begin{grammarentry}{Merksatz}{Warum \glqq gemischt\grqq?}
\GermanText{Diese Nomen verbinden Eigenschaften der schwachen und der starken Deklination: \textbf{-n/-en} in den abhängigen Kasus und zusätzlich ein \textbf{-s} im Genitiv.}
\EnglishText{These nouns combine features of weak and strong declension: \textbf{-n/-en} in the dependent cases and an additional \textbf{-s} in the genitive.}
\end{grammarentry}

\section{Dativ Plural: zusätzliches -n \textnormal{(Dative plural: additional -n)}}
\begin{grammarentry}{Grundregel}{Fast alle Pluralnomen}
\GermanText{Im \textbf{Dativ Plural} bekommt das Nomen normalerweise ein zusätzliches \textbf{-n}, wenn die Pluralform nicht bereits auf \textbf{-n/-en} endet und wenn sie nicht auf \textbf{-s} endet.}
\EnglishText{In the \textbf{dative plural}, the noun normally takes an additional \textbf{-n} if the plural form does not already end in \textbf{-n/-en} and if it does not end in \textbf{-s}.}
\end{grammarentry}

\rowcolors{2}{RowBlueLight}{RowBlueDark}
\begin{longtable}{>{\bfseries}p{3.4cm}|p{3.5cm}|p{3.7cm}|p{4.2cm}}
\rowcolor{HeaderBlueDark}
\color{white}\textbf{Singular (Singular)} &
\color{white}\textbf{Plural (Plural)} &
\color{white}\textbf{Dativ Plural} &
\color{white}\textbf{Regel (Rule)} \\
das Kind & die Kinder & den Kindern & zusätzlich -n \\
das Buch & die Bücher & den Büchern & zusätzlich -n \\
die Frau & die Frauen & den Frauen & -n schon vorhanden \\
das Auto & die Autos & den Autos & kein zusätzliches -n nach -s \\
\end{longtable}

\begin{exbox}
\begin{itemize}
\ExampleItem{Ich spiele mit den Kindern.}{I play with the children.}
\ExampleItem{Die Informationen stehen in den Büchern.}{The information is in the books.}
\ExampleItem{Wir fahren mit den Autos.}{We are driving with the cars.}
\end{itemize}
\end{exbox}

\section{Der Sonderfall \glqq Herz\grqq \textnormal{(The special case \emph{Herz})}}
\begin{grammarentry}{das Herz}{unregelmäßige Mischform}
\GermanText{\textbf{das Herz} ist ein wichtiges neutrales Sonderwort. Im heutigen Standarddeutsch lauten die typischen Formen: \textbf{das Herz, das Herz, dem Herzen, des Herzens}. Der Plural lautet \textbf{die Herzen}.}
\EnglishText{\textbf{das Herz} is an important neuter special case. In present-day Standard German, the typical forms are: \textbf{das Herz, das Herz, dem Herzen, des Herzens}. The plural is \textbf{die Herzen}.}
\end{grammarentry}

\rowcolors{2}{RowBlueLight}{RowBlueDark}
\begin{longtable}{>{\bfseries}p{4cm}|p{7cm}}
\rowcolor{HeaderBlueDark}
\color{white}\textbf{Kasus (Case)} &
\color{white}\textbf{Form (Form)} \\
Nominativ Singular & das Herz \\
Akkusativ Singular & das Herz \\
Dativ Singular & dem Herzen \\
Genitiv Singular & des Herzens \\
Plural & die Herzen \\
\end{longtable}

\section{Eigennamen im Genitiv \textnormal{(Proper names in the genitive)}}
\begin{grammarentry}{Personennamen}{meist -s ohne Artikel}
\GermanText{Personennamen ohne Artikel bekommen im Genitiv normalerweise \textbf{-s}: \textbf{Annas Buch, Peters Auto, Goethes Werke}.}
\EnglishText{Personal names used without an article normally take \textbf{-s} in the genitive: \emph{Annas Buch, Peters Auto, Goethes Werke}.}

\GermanText{Endet ein Name bereits auf einen s-Laut wie \textbf{-s, -ß, -x, -z}, setzt man im Deutschen normalerweise einen \textbf{Apostroph ohne zusätzliches s}: \textbf{Max' Auto, Hans' Bruder}.}
\EnglishText{If a name already ends in an s-sound such as \textbf{-s, -ß, -x, -z}, German normally uses an \textbf{apostrophe without an additional s}: \emph{Max' Auto, Hans' Bruder}.}
\end{grammarentry}

\section{Wann das Nomen keine Kasusendung bekommt \textnormal{(When the noun gets no case ending)}}
\begin{grammarentry}{Sehr häufig keine Veränderung}{Normale feminine Nomen und viele andere Formen}
\GermanText{Viele Nomen verändern ihre Form in den meisten Kasus überhaupt nicht. Besonders feminine Nomen im Singular bleiben normalerweise gleich: \textbf{die Frau, der Frau, der Frau}. Der Kasus wird dann vor allem durch Artikel, Adjektiv oder Satzfunktion gezeigt.}
\EnglishText{Many nouns do not change their form in most cases at all. Feminine nouns in the singular, in particular, normally remain unchanged: \emph{die Frau, der Frau, der Frau}. The case is then shown mainly by the article, adjective, or sentence function.}
\end{grammarentry}

\begin{exbox}
\begin{itemize}
\ExampleItem{Ich sehe die Frau.}{I see the woman.}
\ExampleItem{Ich helfe der Frau.}{I help the woman.}
\ExampleItem{Das Auto der Frau ist neu.}{The woman's car is new.}
\end{itemize}
\end{exbox}

\section{Nomenendung oder Adjektivendung? \textnormal{(Noun ending or adjective ending?)}}
\begin{grammarentry}{Nicht verwechseln}{Zwei verschiedene Systeme}
\GermanText{In \textbf{mit einem kleinen Kind} ist \textbf{-en} in \glqq kleinen\grqq eine \textbf{Adjektivendung}. Das Nomen \glqq Kind\grqq hat im Dativ Singular keine zusätzliche Endung.}
\EnglishText{In \emph{mit einem kleinen Kind}, the \textbf{-en} in \emph{kleinen} is an \textbf{adjective ending}. The noun \emph{Kind} has no additional ending in the dative singular.}

\GermanText{In \textbf{mit kleinen Kindern} ist dagegen das letzte \textbf{-n} in \glqq Kindern\grqq eine echte \textbf{Dativ-Plural-Endung des Nomens}.}
\EnglishText{In \emph{mit kleinen Kindern}, however, the final \textbf{-n} in \emph{Kindern} is a genuine \textbf{dative plural ending of the noun}.}
\end{grammarentry}

\section{Ein praktischer Entscheidungsweg \textnormal{(A practical decision path)}}
\begin{grammarentry}{Schritt für Schritt}{Welche Endung brauche ich?}
\GermanText{\textbf{1. Ist das Nomen Plural und steht es im Dativ?} Dann prüfe das zusätzliche \textbf{-n}: \textbf{mit den Kindern}, aber \textbf{mit den Frauen}, \textbf{mit den Autos}.}
\EnglishText{\textbf{1. Is the noun plural and in the dative?} Then check the additional \textbf{-n}: \emph{mit den Kindern}, but \emph{mit den Frauen}, \emph{mit den Autos}.}

\GermanText{\textbf{2. Ist es ein gemischtes Nomen wie \glqq Name\grqq?} Dann: \textbf{den Namen, dem Namen, des Namens}.}
\EnglishText{\textbf{2. Is it a mixed-declension noun such as \emph{Name}?} Then use: \emph{den Namen, dem Namen, des Namens}.}

\GermanText{\textbf{3. Ist es ein normales maskulines oder neutrales Nomen im Genitiv Singular?} Dann normalerweise \textbf{-s/-es}: \textbf{des Lehrers, des Kindes}.}
\EnglishText{\textbf{3. Is it a normal masculine or neuter noun in the genitive singular?} Then normally use \textbf{-s/-es}: \emph{des Lehrers, des Kindes}.}

\GermanText{\textbf{4. Trifft keiner dieser Fälle zu?} Dann bekommt das Nomen wegen des Kasus meistens keine eigene Endung.}
\EnglishText{\textbf{4. Does none of these cases apply?} Then the noun usually does not get its own ending because of case.}
\end{grammarentry}

\section{Typische Fehler \textnormal{(Typical mistakes)}}


\begin{grammarentry}{Fehler 1}{Dativ-Plural-n vergessen}
\GermanText{Falsch: \textbf{mit den Kinder}. Richtig: \textbf{mit den Kindern}.}
\EnglishText{Wrong: \emph{mit den Kinder}. Correct: \emph{mit den Kindern}.}
\end{grammarentry}

\begin{grammarentry}{Fehler 2}{Genitivendung vergessen}
\GermanText{Falsch: \textbf{die Farbe des Auto}. Richtig: \textbf{die Farbe des Autos}.}
\EnglishText{Wrong: \emph{die Farbe des Auto}. Correct: \emph{die Farbe des Autos}.}
\end{grammarentry}

\begin{grammarentry}{Fehler 3}{-s bei femininen Nomen hinzufügen}
\GermanText{Falsch: \textbf{die Tasche der Fraus}. Richtig: \textbf{die Tasche der Frau}.}
\EnglishText{Wrong: \emph{die Tasche der Fraus}. Correct: \emph{die Tasche der Frau}.}
\end{grammarentry}

\begin{grammarentry}{Fehler 4}{Pluralbildung und Kasusendung verwechseln}
\GermanText{In \textbf{die Kinder} gehört \textbf{-er} zur Pluralform des Wortes. Erst in \textbf{mit den Kindern} kommt wegen des Dativs zusätzlich \textbf{-n} dazu.}
\EnglishText{In \emph{die Kinder}, \textbf{-er} belongs to the plural form of the word. Only in \emph{mit den Kindern} is an additional \textbf{-n} added because of the dative.}
\end{grammarentry}

\section{Pluralbildung ist etwas anderes \textnormal{(Plural formation is a different issue)}}

\begin{grammarentry}{Kasus oder Numerus?}{Case or number?}
\GermanText{Wenn aus Singular Plural wird, verändert sich das Nomen oft ebenfalls: \textbf{Stadt $\rightarrow$ Städte}, \textbf{Kind $\rightarrow$ Kinder}, \textbf{Frau $\rightarrow$ Frauen}. Das ist aber \textbf{Pluralbildung} und nicht eine Veränderung wegen Nominativ, Akkusativ, Dativ oder Genitiv.}
\EnglishText{When singular becomes plural, the noun often changes as well: \textbf{Stadt $\rightarrow$ Städte}, \textbf{Kind $\rightarrow$ Kinder}, \textbf{Frau $\rightarrow$ Frauen}. However, this is \textbf{plural formation}, not a change caused by nominative, accusative, dative, or genitive.}

\GermanText{Danach kann zusätzlich noch eine Kasusendung dazukommen: \textbf{Stadt $\rightarrow$ Städte $\rightarrow$ Städten}.}
\EnglishText{After that, an additional case ending may still be added: \textbf{Stadt $\rightarrow$ Städte $\rightarrow$ Städten}.}
\end{grammarentry}

\section{Übersichtstabelle \textnormal{(Summary table)}}

\rowcolors{2}{RowBlueLight}{RowBlueDark}
\begin{longtable}{>{\bfseries}p{3.3cm}|p{4.0cm}|p{4.0cm}|p{3.0cm}}
\rowcolor{HeaderBlueDark}
\color{white}\textbf{Situation (Situation)} &
\color{white}\textbf{Änderung (Change)} &
\color{white}\textbf{Beispiel (Example)} &
\color{white}\textbf{Wichtig (Important)} \\
Dativ Plural & meistens + -n & Städte $\rightarrow$ Städten & sehr wichtig \\
Genitiv Singular Mask./Neutr. & + -s oder -es & Haus $\rightarrow$ Hauses & sehr wichtig \\
Gemischte Deklination & -n/-en; Genitiv zusätzlich -s & Name $\rightarrow$ Namen/Namens & wichtig \\
Sonderfall Herz & Dat. Herzen; Gen. Herzens & Herz $\rightarrow$ Herzen/Herzens & merken \\
Eigennamen im Genitiv & meistens + -s & Peter $\rightarrow$ Peters & häufig \\
Pluralbildung & verschiedene Endungen/Umlaut & Stadt $\rightarrow$ Städte & kein Kasus \\
\end{longtable}


\section{Merksystem für B1 \textnormal{(Memory system for B1)}}

\begin{grammarentry}{Drei Fragen}{Three questions}
\GermanText{Wenn du wissen willst, ob sich das Ende eines Nomens ändern muss, prüfe in dieser Reihenfolge:}
\EnglishText{If you want to know whether the ending of a noun must change, check in this order:}

\GermanText{1. Ist es \textbf{Dativ Plural}? $\rightarrow$ meistens \textbf{-n}.}
\EnglishText{1. Is it \textbf{dative plural}? $\rightarrow$ usually \textbf{-n}.}

\GermanText{2. Ist es \textbf{Genitiv Singular Maskulinum oder Neutrum}? $\rightarrow$ meistens \textbf{-s/-es}.}
\EnglishText{2. Is it \textbf{genitive singular masculine or neuter}? $\rightarrow$ usually \textbf{-s/-es}.}

\GermanText{3. Ist es ein \textbf{Sonderfall} wie \textbf{Name} oder \textbf{Herz}? $\rightarrow$ eigene Formen lernen.}
\EnglishText{3. Is it a \textbf{special case} such as \textbf{Name} or \textbf{Herz}? $\rightarrow$ learn its individual forms.}
\end{grammarentry}

\section{Zusammenfassung \textnormal{(Summary)}}
\begin{grammarentry}{Die drei wichtigsten Regeln}{Merkschema}
\GermanText{\textbf{1. Genitiv Singular, maskulin/neutrum:} meistens \textbf{-s/-es}.}
\EnglishText{\textbf{1. Genitive singular, masculine/neuter:} usually \textbf{-s/-es}.}

\GermanText{\textbf{2. Dativ Plural:} normalerweise zusätzlich \textbf{-n}, außer die Pluralform endet bereits auf \textbf{-n/-en} oder auf \textbf{-s}.}
\EnglishText{\textbf{2. Dative plural:} normally an additional \textbf{-n}, unless the plural already ends in \textbf{-n/-en} or in \textbf{-s}.}

\GermanText{\textbf{3. Sonder- und Mischformen:} besonders \textbf{Name $\rightarrow$ Namens} und \textbf{Herz $\rightarrow$ Herzen/Herzens} bewusst lernen.}
\EnglishText{\textbf{3. Special and mixed forms:} especially learn \textbf{Name $\rightarrow$ Namens} and \textbf{Herz $\rightarrow$ Herzen/Herzens} consciously.}
\end{grammarentry}



\end{document}
